% ============================================================
% CHAPTER: METHODOLOGY
% ============================================================

\chapter{Methodology}

The overall approach for this project focuses on leveraging vocal characteristics for stress prediction. The strategy integrates data analytics and theoretical models to detect emotion and depression levels effectively. A Convolutional Neural Network (CNN) model is utilized, trained on the RAVDESS dataset (Ryerson Audio-Visual Database of Emotional Speech and Song) and the DAIC-WOZ dataset for depression analysis.

\section{Emotion Analysis}

\begin{enumerate}
    \item \textbf{Input Audio Data:} Begin by gathering raw audio files from the RAVDESS dataset, which includes labeled emotional speech and song recordings.
    \item \textbf{Data Extraction:}
    \begin{itemize}
        \item Extract metadata such as labels and identifiers from the audio file names to organize the dataset.
        \item Create an initial DataFrame to serve as the foundation for further processing.
    \end{itemize}
    \item \textbf{Data Augmentation:}
    \begin{itemize}
        \item Introduce diversity in the dataset by adding noise and altering the pitch of the audio files.
        \item Update the DataFrame with these augmented samples to expand the training set.
    \end{itemize}
    \item \textbf{Feature Extraction:}
    \begin{itemize}
        \item Extract Mel-frequency Cepstral Coefficients (MFCC) from the audio data to capture relevant features.
        \item Build a final DataFrame containing these extracted features for model training.
    \end{itemize}
    \item \textbf{Data Transformation:}
    \begin{itemize}
        \item Transform class labels into a non-hot encoded format to prepare them for CNN model training.
        \item Ensure the dataset is structured and ready for machine learning tasks.
    \end{itemize}
    \item \textbf{Model Training:}
    \begin{itemize}
        \item Train the CNN model using the processed and transformed data.
        \item Optimize the model by fine-tuning hyperparameters and evaluating performance.
    \end{itemize}
\end{enumerate}

\section{Depression Analysis}

\subsection{Database Utilization}
The DAIC-WOZ database is used for predicting user depression levels based on vocal features extracted from audio files.

\subsection{Data Preprocessing}
\begin{itemize}
    \item \textbf{Audio File Segmentation:}
    \begin{itemize}
        \item The audio data is split into two parts: the speaker's segment and the participant's segment.
        \item The participant's audio segment is further divided into 2-minute duration chunks to standardize the input length for analysis.
    \end{itemize}
    \item \textbf{Noise and Pitch Augmentation:}
    \begin{itemize}
        \item Noise and pitch variations are added to the audio chunks to prevent overfitting and improve model performance under varying noise conditions.
    \end{itemize}
\end{itemize}

\subsection{Feature Extraction}
The preprocessed audio chunks are converted into MFCC (Mel-frequency Cepstral Coefficients) features. These features serve as inputs for predicting stress or depression levels, as they effectively capture vocal characteristics.

\subsection{Model Selection and Architecture}
A Convolutional Neural Network (CNN) model is implemented for prediction.
\begin{itemize}
    \item \textbf{Input Layer:} Accepts MFCC features with a specified input shape.
    \item \textbf{Convolutional Layers:} Multiple Conv1D layers are used, each with different filter sizes and padding configurations, followed by ReLU activation.
    \item \textbf{Batch Normalization:} Ensures stable and faster training by normalizing the outputs of convolutional layers.
    \item \textbf{Dropout Layers:} Prevent overfitting by randomly disabling a fraction of neurons during training.
    \item \textbf{MaxPooling Layers:} Reduce dimensionality while retaining key information.
    \item \textbf{Fully Connected Layer:} Maps the flattened features to the output classes.
    \item \textbf{Output Layer:} Uses a softmax activation function to output probabilities for each target class.
\end{itemize}

\section{Website Development Methodology}

\begin{enumerate}
    \item \textbf{Frontend Development:}
    \begin{itemize}
        \item The website is developed using HTML, CSS, and JavaScript to provide an interactive and user-friendly interface.
        \item The frontend collects responses from users through a structured questionnaire containing 40 questions derived from the DAIC-WOZ dataset.
    \end{itemize}
    \item \textbf{Backend Development:}
    \begin{itemize}
        \item Python Flask is used for the backend to handle data processing and communication between the frontend and server.
        \item User responses are saved in audio format for further analysis.
    \end{itemize}
    \item \textbf{Audio Data Preprocessing:}
    \begin{itemize}
        \item The audio files are split into two durations:
        \begin{itemize}
            \item 3-second segments: Used to predict emotion types like calm, angry, disgust, fearful, and others.
            \item 2-minute chunks: Used for stress and depression level detection.
        \end{itemize}
    \end{itemize}
    \item \textbf{Emotion Detection:}
    \begin{itemize}
        \item A CNN model is employed to classify emotions based on the 3-second audio segments.
        \item The output categorizes emotions into predefined types, aiding in understanding the emotional state of the user.
    \end{itemize}
    \item \textbf{Depression Detection:}
    \begin{itemize}
        \item Another CNN model is used to predict stress levels based on the 2-minute audio chunks.
        \item The model leverages vocal features extracted from the audio to determine the user's depression level.
    \end{itemize}
    \item \textbf{Data Analysis and Diagnosis:}
    \begin{itemize}
        \item Extracted data, including emotions and stress levels, is further analyzed to derive meaningful insights.
        \item These insights are used to perform comprehensive diagnostics and provide personalized feedback or recommendations.
    \end{itemize}
    \item \textbf{Integration and Workflow:}
    \begin{itemize}
        \item The frontend, backend, and models are seamlessly integrated to enable a smooth user experience.
        \item The system ensures accurate prediction and robust analysis of emotional and mental health data.
    \end{itemize}
\end{enumerate}
